% Subsection: why CGS's gains are smaller on the random DAGs.
% Every number is a macro generated by paper/why_random.py (outputs/tables/why_random_macros.tex) and
% checked by paper/verify.py; the table body is outputs/tables/tab_why_random_rows.tex.
% Copy outputs/tables/ next to the manuscript as tables/ (as for the other tables), then \input this file.
\input{tables/why_random_macros.tex}

\subsection{Why the gains are smaller on random DAGs}
\label{sec:why-random}

CGS is shorter than every baseline on average on the random DAGs as well, but by less than on the
workflows: \WRRandVsAeft\,\% against AEFT and \WRRandVsHeft\,\% against HEFT, compared with \WRWfVsAeft\,\%
and \WRWfVsHeft\,\% on the workflows. The difference is concentrated at low communication. At CCR~10 the
gain over AEFT (\WRRHighVsAeft\,\%) exceeds that on the workflows, while at CCR~0.1 it is \WRRLowVsAeft\,\%.
Table~\ref{tab:why-random} shows that the random graphs offer CGS's two mechanisms little to act on, and
that at CCR~0.1 there is little left for any heuristic to gain.

\paragraph{The joins CGS targets are rare.}
The generator gives each task a number of parents equal to half the width of the previous level, drawn from
the two preceding levels, so the random DAGs are dense: a non-entry task has \WRRandParents{} parents on average, against
\WRWfParents{} in the workflows, and only \WRRandFewParents\,\% of tasks have one or two parents, against
\WRWfFewParents\,\% (Table~\ref{tab:why-random}). Joint placement applies only to a join whose remaining
parents number at most two, so on a random DAG it acts on the last two of about \WRRandJoinParents{} parents.
By then the join's ready time is usually fixed by the parents placed earlier. Instrumenting every closure on
the test split shows that this holds for \WRRandSetEarlier\,\% of closures on the random DAGs, against
\WRWfSetEarlier\,\% on the workflows (\WRWfMinSetEarlier--\WRWfMaxSetEarlier\,\% by family). The part of the
join's ready time that the closure's parents add beyond what the earlier parents impose is \WRRandShift\,\% on
average, against \WRWfShift\,\% (\WRWfMinShift--\WRWfMaxShift\,\% by family). Closures still change
\WRRandChanged\,\% of the random schedules, but at low communication they shorten about as many as they
lengthen (\WRRLowShorter{} against \WRRLowLonger{} at CCR~0.1, \WRRMidShorter{} against \WRRMidLonger{} at
CCR~1, compared with \WRWfShorter{} against \WRWfLonger{} on the workflows), so joint placement contributes
\WRRLowJoint\,\% at CCR~0.1 and \WRRMidJoint\,\% at CCR~1. Only at CCR~10, where the communication on the last parents' edges is large, does
it contribute (\WRRHighJoint\,\%, against \WRWfJoint\,\% on the workflows).

\paragraph{The downstream table is less informative.}
The table $T$ that CGS uses in processor selection, like AEFT's IOCT, is optimistic: it assumes that a child
can start once the data of the task being placed arrives. With \WRRandParents{} parents per task that
assumption rarely holds, so $T$ separates processors poorly. Removing $T$ from selection costs CGS
\WRWfTable\,\% on the workflows (\WRWfMinTable--\WRWfMaxTable\,\% by family) but only \WRRLowTable\,\% and
\WRRMidTable\,\% on the random DAGs at CCR~0.1 and~1, and \WRRHighTable\,\% at CCR~10.

\paragraph{At low communication there is little left to gain.}
Let the bound of an instance be the larger of the critical path with every task at its minimum cost and the
total minimum cost divided by the number of processors. It ignores communication, so it is close to tight
only at CCR~0.1, where AEFT's makespan is a median \WRRLowAeftGap\,\% above it and the best of the seven
baselines \WRRLowBestGap\,\% above it: no method can gain more than that. The strongest baselines also agree:
the median spread among AEFT, PMMS and HEFT is \WRRLowSpread\,\% at CCR~0.1 and \WRRMidSpread\,\% at CCR~1,
against \WRWfSpread\,\% on the workflows. When good heuristics already produce nearly the same schedule, a gain
of \WRRLowVsAeft\,\% over AEFT is what remains.

\paragraph{Summary.}
CGS's advantage comes from joins with few parents and from communication large enough for their placement to
matter. The workflows have both; the dense random DAGs have neither at CCR~0.1 and~1, and at CCR~0.1 they leave
almost no room for improvement. The smaller gains on the random DAGs therefore reflect the graphs rather than a
weakness of the method: where its mechanisms do not apply, CGS still matches or beats every baseline on average.

Two limits apply. The random DAGs come from one generator setting (edge density~0.5); the explanation predicts
larger gains on sparser random graphs, which we have not tested. The bound ignores communication, so the
headroom argument is made for CCR~0.1 only.

\begin{table*}[t]
\centering
\caption{Why CGS gains less on the random DAGs (test split). Parents: mean number of parents of a non-entry
task. 1--2 parents: share of non-entry tasks with one or two parents. Set earlier: share of closures whose
join ready time was already set by parents placed before the closure. Added: mean share of the join's ready time
that the closure's parents add beyond what the earlier parents impose. Joint and Table: mean paired gain (\%) of CGS over the variant without
joint placement and without $T$ in processor selection. Spread: median relative difference between the best
and worst of AEFT, PMMS and HEFT. vs AEFT: mean paired gain of CGS over AEFT. All columns after Parents are in \%.}
\label{tab:why-random}
\footnotesize
\setlength{\tabcolsep}{4.5pt}
\begin{tabular}{lrrrrrrrrr}
\toprule
 & & \multicolumn{2}{c}{Structure} & \multicolumn{2}{c}{Closures} & \multicolumn{2}{c}{Component gain} & & \\
\cmidrule(lr){3-4}\cmidrule(lr){5-6}\cmidrule(lr){7-8}
Instances & $N$ & Parents & 1--2 parents & Set earlier & Added & Joint & Table & Spread & vs AEFT \\
\midrule
\csname @@input\endcsname tables/tab_why_random_rows.tex % primitive input: \input is not expandable inside tabular
\bottomrule
\end{tabular}
\end{table*}
